%% kapa: exam=1 round=24 year=2013 period=1 form=A subject=civil_law subject_ko=민법 seq=2 number=2
%% source: Q-Net 감정평가사 기출문제 (cst00309) artlSeq=5073802 file=extracted/1008/5073802/2013년 제24회 감정평가사 1차시험 기출문제/1교시(A형).hwp
%% conversion: kordoc 4.24.51 -> tools/gen1.py (auto); review=proofread-sonnet; flags=none
\begin{question}{2}
실종선고와 그 취소에 관한 설명으로 옳은 것은? (다툼이 있으면 판례에 의함)
\begin{choices}
  \item 실종선고 또는 인정사망이 있는 경우 사망한 것으로 본다.
  \item 실종선고의 취소는 실종기간 경과 후 취소 전에 선의로 한 행위의 효력에 영향을 미치지 않는다.
  \item 실종선고를 받은 자가 생환한 경우 그 실종선고는 효력을 잃는다.
  \item 실종선고가 취소되면 실종선고를 원인으로 재산을 취득한 자는 받은 이익의 전부를 반환하여야 한다.
  \item 동일인에 대하여 2차례의 실종선고가 있는 경우 첫번째 실종선고를 기준으로 상속관계를 판단한다.
\end{choices}
\end{question}
\answer{5}
